\centerline{\bf \Huge Антигеомный разнобой}
\centerline{\bf \Huge Решения}

\problem{1}
Рассмотрим график приведенного квадратного трехчлена с положительным
дискриминантом. Проведем окружность через три точки пересечения этого
графика с осями координат. Докажите, что все такие окружности проходят
через фиксированную точку на плоскости.

\textbf{Решение.}
Докажем, что все такие окружности проходят через точку $(0,1)$.
Обозначим через $x_1$ и $x_2$ абсциссы точек пересечения графика
произвольного квадратного трехчлена с осью $Ox$, а через $q$ — ординату
точки пересечения с осью $Oy$. Заметим, что числа $x_1$ и $x_2$ — корни
квадратного трехчлена со свободным коэффициентом $q$, поэтому по теореме
Виета $x_1x_2=q$. Однако левая часть равенства — это в точности степень
точки $O$ (начала координат) относительно окружности, проходящей через
точки $x_1$, $x_2$, $q$. С другой стороны, $1\cdot q=x_1x_2$, поэтому
точка с ординатой 1 на оси $Oy$ также лежит на этой окружности, что и
требовалось.

\problem{2}
 языке племени УЫ всего две буквы: У и Ы. Словом считается любая
последовательность из $2 n$ букв У и $2 n$ букв Ы (число $n$ дано и
фиксировано). Язьковеды называют слова похожими, если одно можно
получить из другого одной перестановкой двух соседних букв У и Ы. Какое
наибольшее количество слов можно выписать на доску так, чтобы любые два
из выписанных слов не были похожи?

\textbf{Решение.}
Возьмем все слова, которых $C_{4n}^{2n}$. Далее, разобьем слова на блоки
по две подряд идущие буквы. Если в слове нашелся блок, который состоит
из разных букв, сопоставим ему похожее слово, которое получается
перестановкой букв в самом левом таком блоке. Из каждой пары мы можем
взять не более одного слова, а также мы можем взять все слова, которые
состоят из блоков УУ и ЫЫ. Последних ровно $C_{2n}^n$, а пар похожих —
$(C_{4n}^{2n}-C_{2n}^n)/2$. Значит, всего можно взять не более
$(C_{4n}^{2n}-C_{2n}^n)/2+C_{2n}^n=(C_{4n}^{2n}+C_{2n}^n)/2$ слов.

\problem{3}
Пусть $a$ и $b$ — натуральные числа, такие, что число $a+b^3$ делится на
число $a^2+3ab+3b^2-1$. Докажите, что число $a^2+3ab+3b^2-1$ делится на
точный куб натурального числа, большего, чем 1.

Пусть $a^2+3ab+3b^2-1=n$. Тогда $a+b^3=nx$ для некоторого натурального
$x$. Заметим, что
$$(a+b)^3=a^3+3a^2b+3ab^2+b^3=a(a^2+3ab+3b^2-1)+(a+b^3)=an+xn,$$ поэтому
$(a+b)^3$ делится на $n$. Предположим, что каждый простой делитель,
входящий в $n$, входит в него в степени не выше 2. Тогда этот же простой
делитель должен входить в число $a+b$, а потому
$(a+b)^2\geq n=a^2+3ab+3b^2-1$. Но это неверно, т.к.
$$a^2+3ab+3b^2-1=(a+b)^2+b^2+(b^2-1)> (a+b)^2.$$ Значит, $n$ делится на
куб какого-то простого числа, что и требовалось доказать.

\textbf{Решение.}

\problem{4}
Дано натуральное $n$. В ряд записаны $n$ целых чисел. Два человека
играют в игру. За один шаг первый указывает на несколько из них,
записанных подряд, а второй либо увеличивает каждое из указанных чисел
на 1, либо уменьшает каждое из них на 1. Найдите наибольшее $k$, такое,
для которого первый всегда за несколько шагов сможет добиться, чтобы
хотя бы $k$ чисел стали делиться на 3.

\textbf{Решение.}
Очевидно, что второй всегда сможет составить одно из чисел не делящимся
на 3, поэтому $k<n$. Докажем, что $k=n-1$. Будем вести доказательство по
индукции. При $n=2$ утверждение очевидно. Для произвольного $n>2$
рассмотрим два числа, которые стоят по краям отрезка. Если одно из них
делится на 3, то будем рассматривать отрезок без этого числа и применим
предположение индукции. В противном случае либо у этих чисел разные
ненулевые остатки по модулю 3, либо одинаковые. Если остатки одинаковые,
то попросим поменять одно из этих чисел. Тогда оно либо станет делиться
на 3, либо будет иметь другой ненулевой остаток, что сводит нашу задач к
случаю разных остатков на концах отрезка. Но тогда, попросив поменять
числа во всем отрезке, мы получим, что одно из чисел на конце отрезка
станет делиться на 3.

\problem{5}
Найти наибольшее натуральное $n$, при котором сумма
$[\sqrt{1}]+[\sqrt{2}]+\ldots+[\sqrt{n}]$ является простым числом.

\textbf{Решение.}
Выберем произвольное натуральное число $n$ и рассмотрим такое
натуральное $k$, что $k^2\leq n\leq k^2+2k$. Для всех чисел $m$ из
отрезка $[k^2, k^2+2k]$ длины $2k+1$ имеем: $[\sqrt{m}]=k$. Таким
образом, сумма $[\sqrt{1}]+[\sqrt{2}]+\ldots+[\sqrt{n}]$ устроена
следующим образом. Пусть $n=k^2+\ell$, где $0\leq\ell\leq 2k$. Тогда в
нашей сумме есть слагаемое 1, которое повторяется 3 раза, есть слагаемое
2, которое повторяется 5 раз, …, слагаемое $i$, которое повторяется
$2i+1$ раз, …, слагаемое $k$, которое повторяется $\ell+1$ раз. Вычислим
эту сумму явно: $$\begin{gathered}
[\sqrt{1}]+[\sqrt{2}]+\ldots+[\sqrt{n}] = \sum\limits_{i=1}^{k-1} i(2i+1)+k(\ell+1) = 2\sum\limits_{i=1}^{k-1} i^2+\sum\limits_{i=1}^{k-1} i + k(\ell+1) = \\ = \frac{2(k-1)k\big(2(k-1)+1\big)}{6}+\frac{(k-1)k}{2}+k(\ell+1) = \frac{(k-1)k(4k+1)}{6}+k(\ell+1).\end{gathered}$$

Заметим, что если $k>6$, то у первого и второго слагаемого есть общий
простой делитель (поскольку в обоих присутствует множитель $k$). Значит,
сумма не может быть простой при $k\geq 7$ и $n\geq k^2=49$. При $n=48$
получаем сумму, равную $203=7\cdot 29$ — не подходит, а при $n=47$ —
сумму, равную простому числу 197. Таким образом, наибольшее $n$ равно
47.

\problem{6}
Дана доска $2m\times 2n$, раскрашенная в шахматном порядке. Сколькими
способами можно поставить на белые клетки $mn$ шашек так, чтобы никакие
две шашки не стояли бы на одной клетке и чтобы никакие две шашки не
располагались бы в клетках, соседних по углу?

\textbf{Решение.}
Разделим таблицу на квадраты $2\times 2$. Ясно, что в каждом квадрате
$2\times 2$ есть ровно одна шашка. Поэтому нам нужно посчитать,
сколькими способами можно расставить по одной шашке в каждый квадрат
$2\times 2$, чтобы клетки с шашками не граничили бы по углу. Без
ограничения общности будем считать, что у всех квадратов $2\times 2$
правая нижняя клетка белая. Будет называть квадрат ПН-квадратом, если
шашка стоит в правой нижней его клетке, и ЛВ-квадратом, если шашка стоит
в левой верхнем углу.

Заметим, что если какой-то квадрат является ПН-квадратом, то и квадраты
слева и сверху от него являются ПН-квадратами. Аналогичное верно и для
ЛВ-квадратов. Таким образом, все ПН-квадраты образуют связную область, и
все ЛВ-квадраты образуют связную область.

Ясно, что количество расстановок шашек зависит от количества способов
разбить наш прямоугольник на две такие области. Линия границы между
этими областями ведет из левого нижнего угла таблицы в правй верхний
угол по линиям сетки. Количество способов провести такую линию равно
соответствующему биномиальному коэффициенту $C_{m+n}^m$.


\problem{7}
а) Алиса и Боб играют в игру. На плоскости отмечены $n$ точек общего
положения, где $n$ — нечетное натуральное число; $n$ — произвольное
натуральное число, большее 1. Алиса и Боб по очереди выбирают пару точек
и соединяют их отрезком. Разрешается соединять только точки, между
которыми еще не проведен отрезок. Проигрывает тот, после чьего хода
образуется цикл нечетной длины. Начинает Алиса. Кто выигрывает при
правильной игре?

б) Что поменяется, если $n$ чётное?

\textbf{Решение (a)}
Докажем, что выиграет Боб. Докажем, что в любой момент времени после
хода Алисы он сможет сделать свой ход и не проиграть. Рассмотрим граф,
возникающий перед ходом Боба. Поскольку он не содержит циклов нечетной
длины, этот граф является двудольным. Пусть $a$ и $b$ — размеры его
долей. Ясно, что проведение любого нового ребра между долями не приведет
к поражению, поэтому, если Боб может провести такое ребро, он его
проведет. Предположим, что Боб не может сделать ход. Тогда перед его
ходом возник полный двудольный граф $K_{a,b}$. Но поскольку $a+b=n$ —
нечетно, числа $a$ и $b$ имеют разную четность, а количество проведенных
ребер равно $ab$ — четно. Однако после хода Алисы проведено нечетное
число ребер. Значит, после ее хода никак не может получиться полный
двудольный граф. Поэтому Боб всегда сможет сделать ход. Таким образом,
он выиграет.

\textbf{Решение (б)}

Докажем, что при $n\equiv_4 2$ Алиса выиграет.
Первым ходом Алиса соединяет любые две вершины. Теперь перед каждым
ходом Боба есть вершины (их множество обозначим через $A$), из которых
выходит хотя бы одно ребро, и вершины (их множество обозначим через
$B$), из которых ребер пока что не выходит. До тех пор, пока $|B|>0$,
будем действовать так:

-   если Боб соединяет две вершины из $A$, Алиса соединит две вершины из
    $B$;

-   если Боб соединит две вершины $b_1$ и $b_2$ из $B$, Алиса соединит
    вершину $b_1$ с любой вершиной из $A$;

-   если Боб соединит вершину из $A$ и вершину $b$ из $B$, то Алиса
    соединит $b$ с любой другой вершиной из $B$.

Поясним, почему возможен последний тип хода. Дело в том, что если
$n=|A|+|B|$ четно, то $|A|$ и $|B|$ имеют одинаковую четность. В
стартовый момент $|A|$ и $|B|$ четны, а после ходов первых двух типов
четность чисел $|A|$ и $|B|$ не меняется ($|A|$ увеличивается на 2, а
$|B|$ уменьшается на 2). Значит, $|B|$ всегда четно, и третий тип хода
также возможен (он также сохраняет четности чисел $|A|$ и $|B|$).

Алиса будет действовать по этому алгоритму до тех пор, пока не останется
изолированных вершин. Когда это произойдет, Алиса будет просто делать
любые допустимые ходы. Предположим, что в какой-то момент времени Алиса
не сможет сделать допустимый ход. Поскольку в получившемся графе нет
нечетных циклов, это полный двудольный граф. Докажем, что это граф
$K_{n/2, n/2}$, Действительно, заметим, что в процессе построения графа
после каждого хода Алисы в графе существует паросочетание на множестве
вершин $A$. Значит, в тот момент, когда $|A|$ станет равным $n$, на
получившемся двудольном графе будет существовать совершенное
паросочетание. Это и означает, что граф состоит из долей одинакового
размера.

Итак, после некоторого хода Боба возникнет граф $K_{n/2,n/2}$,
соедржащий $n^2/4$ ребер. Но $n^2/4$ нечетно, а после хода Боба
проведено четное количество ребер — противоречие. Значит, у Алисы всегда
будет допустимый ход, и она выиграет.

Если же $n\equiv_4 0$, то после первого хода Алисы Боб начнет
использовать ее стратегию и сумеет победить. Таким образом, Алисы сможет
выиграть, только если $n\equiv_4 2$.

При $n\equiv_4 2$ выигрывает Алиса, при всех остальных $n$ — Боб.

\problem{8}
Дан многочлен $P(x)=ax^3+bx^2+cx+d$ с целыми коэффициентами, такими, что
$a\neq 0$. Известно, что существует бесконечно много пар целых чисел
$(x,y)$, для которых справедливо равенство $xP(x)=yP(y)$. Докажите, что
у многочлена $P(x)$ есть целый корень.

\textbf{Решение.}
Преобразуем разность $xP(x)-yP(y)$: $$\begin{gathered}
a(x^4-y^4)+b(x^3-y^3)+c(x^2-y^2)+d(x-y)=0\quad\Longrightarrow \\ \Longrightarrow\quad a(x+y)(x^2+y^2)+b(x^2+xy+y^2)+c(x+y)+d=0.    \end{gathered}$$
Положим $x+y=u$ и $x^2+y^2=v$, тогда $(2au+b)v=-(bu^2+2cu+2d)$. Заметим,
что $v\geq u^2/2$, откуда $u^2|2au+b|\leq |bu^2+2cu+2d|$. Ясно, что это
неравенство имеет лишь конечное количество решений, т.е. существует лишь
конечное возможное количество значений суммы $x+y$, удовлетворяющих
соотношению $xP(x)=yP(y)$. Поскольку самих пар $(x,y)$, удовлетворяющих
соотношению $xP(x)=yP(y)$, бесконечно много, получаем, что существует
такая константа $\alpha$, что $x+y=\alpha$ для бесконечно многих пар
$(x,y)$. Но тогда соотношение $xP(x)=(\alpha-x)P(\alpha-x)$ выполняется
для бесконечно многих целых чисел $x$. Это означает, что это равенство
на самом деле выполняется при всех значениях $x$. Если $\alpha\neq 0$,
то полагая $x=\alpha$, получаем, что $P(\alpha)=0$, т.е. $\alpha$ — это
целый корень многочлена $P$. Если же $\alpha=0$, то $P(x)=-P(-x)$, т.е.
многочлен $P$ нечетный и $P(0)=0$.